% 修订：启用 xeCJK 自动换行；保留原题参数与物理关系。
\documentclass[tikz,border=14pt]{standalone}
\usepackage{fontspec}
\usepackage{xeCJK}
\usepackage{amsmath,amssymb}
\usepackage{xcolor}
\usepackage{tikz}
\usetikzlibrary{arrows.meta,calc,patterns,positioning,decorations.pathmorphing}
\IfFontExistsTF{Heiti SC}{\IfFontExistsTF{Heiti SC}{\setmainfont{Heiti SC}}{\setmainfont{FandolHei-Regular}}\IfFontExistsTF{Heiti SC}{\setsansfont{Heiti SC}}{\setsansfont{TeX Gyre Heros}}}{\IfFontExistsTF{Songti SC}{\setmainfont{Songti SC}}{\setmainfont{FandolSong-Regular}}\IfFontExistsTF{Songti SC}{\setsansfont{Songti SC}}{\setsansfont{FandolSong-Regular}}}
\definecolor{bluefield}{RGB}{24,76,255}
\definecolor{redfield}{RGB}{232,49,35}
\definecolor{orangefield}{RGB}{224,116,10}
\definecolor{greenfield}{RGB}{18,132,86}
\definecolor{purplefield}{RGB}{112,72,190}
\definecolor{linegray}{RGB}{60,68,78}
\definecolor{softblue}{RGB}{238,243,255}
\definecolor{softgreen}{RGB}{238,249,244}
\definecolor{softorange}{RGB}{255,247,235}
\definecolor{softpurple}{RGB}{247,242,255}
\tikzset{
  arrow/.style={-{Latex[length=3.1mm,width=2.1mm]},line width=1.1pt},
  thinarr/.style={-{Latex[length=2.5mm,width=1.7mm]},line width=.85pt},
  axis/.style={-{Latex[length=2.5mm,width=1.7mm]},line width=.8pt,draw=linegray},
  guide/.style={draw=linegray!55,dashed,line width=.65pt},
  fieldline/.style={-{Latex[length=2.5mm,width=1.7mm]},line width=1pt,draw=bluefield},
  box/.style={draw=linegray,rounded corners=3pt,fill=white,align=center,inner sep=6pt,line width=.75pt},
  bluebox/.style={box,draw=bluefield,fill=softblue},
  greenbox/.style={box,draw=greenfield,fill=softgreen},
  orangebox/.style={box,draw=orangefield,fill=softorange},
  purplebox/.style={box,draw=purplefield,fill=softpurple}
}
\IfFontExistsTF{Heiti SC}{\setCJKmainfont[AutoFakeBold=1.5]{Heiti SC}}{\setCJKmainfont[AutoFakeBold=1.5]{FandolHei-Regular}}
\IfFontExistsTF{Heiti SC}{\setCJKsansfont{Heiti SC}}{\setCJKsansfont{FandolHei-Regular}}
\begin{document}

\begin{tikzpicture}[x=1cm,y=1cm]
\path[use as bounding box] (-10,-4.6) rectangle (10,4.6);
\fill[white] (-10,-4.6) rectangle (10,4.6);
\node[font=\Large\bfseries,text width=18.4cm,align=center] at (0,4.05) {模块一：几何、微分运算与积分定理};
\node[bluebox,text width=3.05cm] (n0) at (-7.6,.55) {\textbf{有向几何}\\$d\mathbf l,d\mathbf S,dV$};
\node[greenbox,text width=3.05cm] (n1) at (-3.8,.55) {\textbf{微分运算}\\$\nabla,\nabla\!\cdot,\nabla\!\times$};
\node[orangebox,text width=3.05cm] (n2) at (0,.55) {\textbf{积分对象}\\环量、通量、体积分};
\node[purplebox,text width=3.05cm] (n3) at (3.8,.55) {\textbf{积分定理}\\高斯与斯托克斯};
\node[bluebox,text width=3.05cm] (n4) at (7.6,.55) {\textbf{奇异源}\\$\delta$ 与挖孔区域};
\draw[arrow,bluefield] (n0)--(n1);
\draw[arrow,greenfield] (n1)--(n2);
\draw[arrow,orangefield] (n2)--(n3);
\draw[arrow,purplefield] (n3)--(n4);
\node[greenbox,text width=15.9cm] at (0,-2.45) {先确认对象和方向，再选择算符；遇到奇点时先保留积分式，最后讨论能否局域化。};
\end{tikzpicture}
\end{document}
